\subsection{北边的警报与南方的准备}
长安在开皇初年听见的并不只是新年号。北面来的是突厥军与北周旧部相连的消息，南面隔着长江的是仍有完整朝廷、州县与水军的陈。朝廷先要把关中的仓、府兵和使者接在一处，才有余地选择南下的时机。沙钵略可汗的压力见于\eventref{ST-582-TURKIC-DEFENSE}{开皇二年的北边危机}；它使隋廷不能把军费、马匹和将领全数移到江淮。分化突厥诸部、修补边防，并不是统一叙事前的一段无关序曲，而是让南征不至于立即失去北门的条件。

这种准备也有文书的一面。开皇律令的修订\eventref{ST-583-KAIHUANG-CODE}{在开皇三年颁行}，给官府一套可援用的刑名与行政语言；它不能证明同一时刻每个州县都按同一尺度办案。北齐、北周以来的吏员、地方豪强与乡里关系不会因一部律令而消失。可是新朝若要派官、征粮、处置逃亡和接收降附者，总得先使命令具有可被辨识的形式。制度在这里不是一张静止的清单，而是调动人与物时不断被解释、延误和争执的工具。

南征的实际路线也说明这个限制。隋军在开皇八年分水陆诸军出动，不能只被写成一支从北向南的箭头。长江两岸的渡口、舟楫、前沿州县和陈朝内部的防守安排，决定军队何时能相会、粮食如何跟上，以及降附能否变成持续的控制。建康在次年正月易手，陈后主被俘，见\eventref{ST-589-CHEN-CONQUEST}{建康易手（589）}；这一结果把宫廷的名义交接与整座江南的接管分开了。对隋军而言，入城是军事行动的完成；对新到的官员、旧陈吏员、寺院和乡里而言，户籍、仓储、治安与差役才刚开始重新安排。

江南的反隋起事在590年出现，正把这个差别暴露出来。我们不能由官修史书所列的叛首与讨平路线，复原每一村的动机；旧朝遗民、赋役压力、地方武装与新旧官吏的冲突不必只有一种来源。但\eventref{ST-590-JIANGNAN-REBELLION}{江南起事}至少迫使读者放弃“589年统一已经结束”的读法。新政权能够在地图上合并州郡，并不意味着它已掌握每一段水路、每一座仓和每一份可征的户籍。统一是一连串接收动作，而这些动作的成本落在不同的人身上。

\subsection{洛阳、水道与被拉长的国家}
大业初年东都的经营改变了国家观看空间的方式。洛阳不是只为皇帝增添一座宫城：它处在关中、河北与江淮之间，能把朝廷、仓储与漕运置于同一套较短的指挥距离中。通济渠在605年的开凿，见\eventref{ST-605-TONGJI-CANAL}{通济渠开凿}，让这一设想有了水面上的路径。河渠遗存、后来的整修记录和正史叙事都提示工程的分段性；它们不能把劳役人数或死亡数替我们精确化。更可确定的是，堤岸、船只、征夫和地方供应把远处的江淮家庭也纳入一项都城工程。

水路带来的并不只是压迫。它让谷物、布帛、官员和消息能在季节允许时移动，也让沿线市场、渡口与仓场变成可以被征调和被争夺的节点。正因为如此，同一条通道在丰年、旱涝、战争或征发加剧时会给不同地方造成完全不同的经验。把运河说成单纯“繁荣工程”或单纯“亡国工程”，都会遮掉这种变化。唐初以后仍可利用的路线说明设施并未随隋亡消失；隋末迅速失控又说明设施不能替地方社会承担维护与供给。

对高句丽的三次远征把这个问题推到更远的东北。612年的战事，见\eventref{ST-612-GOGURYEO-INVASION}{高句丽远征}，要求隋廷同时组织北方兵员、陆路转输、水军与辽东行进。汉文史书叙述多以隋廷决策和战果组织材料，高句丽方面留下的可互证材料较少，故而不宜把行军人数、每日路线或对方社会的反应写成无疑的数字。可是辽东的距离、河流与城防本身已足以说明：中央能够下令，并不等于粮秣和人力会按命令抵达。

当征发、战败、地方军政网络与继承不安交错时，原先被用来缩短距离的道路也让消息和兵员更快卷入各地冲突。李渊在617年由太原起兵，见\eventref{ST-617-LI-YUAN-TAIYUAN}{太原起兵}，接住的正是尚未崩毁的一部分关中军政资源、交通和名义。隋的结束因而不该被缩成一位皇帝的性格故事；它更像一个被拉得过长的调度网络，在边疆、都城、水路和地方同时失去可以承重的接点。

\begin{textwindow}{“开皇”与“大业”之间不能只剩一条道德曲线}
《隋书》由唐初官修，其帝纪将开皇与大业安排进后来已知的兴亡结局。《资治通鉴》又把同一材料放进治乱的连续判断。它们能够帮助我们确定改元、工程与战争的先后，却不能让后来的结局代替当时州县、军人和行旅所拥有的信息。读者在两种年号之间应追问道路、仓储、边防与征役怎样相互牵连，而不把“盛”与“亡”当作已经解释一切的标签。
\end{textwindow}

\subsection{关中朝廷怎样把一纸受禅变成可以行走的命令}
受禅礼是在宫城内完成的，开皇元年的难处却在宫城外展开。北周末年的官署、军府、仓廪和驿道没有在改国号的一刻变成同一种东西。杨坚的朝廷首先面对的是一群必须继续办事的人：掌文书的吏员要知道新年号从哪一天写起，守仓者要知道粮帛向何处开支，北面戍守的将领要判断关中政局是否会使补给中断，原来依北周官爵和婚姻网络立足的家族则要衡量留任、观望或另结盟的代价。\eventanchor{ST-581-SUI-ACCESSION}{隋受禅（581）}所能确定的，是政权在二月完成了名义转换；它真正要取得的，是这些分散节点继续承认命令能够抵达、供给能够兑现。

《隋书·高祖纪上》把受禅写成一组诏命和仪式，《北史》与《资治通鉴》可互校其月日。这样的材料很适合观察朝廷如何宣布新秩序，却不适合替每一处州县写出即时的欢呼或恐惧。尤其不能把唐初史官后来看到的统一，倒投为关中人在581年已经拥有的把握。对当时的官府而言，北周静帝退位只处理了最高名义；旧有地方官是否守住治所、军府能否抽调、道路上是谁在传递文书，都还要通过一连串小而具体的确认来完成。

受禅后的政治并不是把旧朝人物一概逐走。新朝需要熟悉户籍、仓储、礼制与边防的人，也需要让留下的人明白，原有关系不能再单独决定资源的流向。任命、赦令和礼仪因此既是宣告，也是试探：哪一处能按时回报，哪一支部队仍听调遣，哪一位地方豪强借新旧交替扩大自己的空间。史书保存的是中枢发出的结论，地方的迟疑常只有在后来叛乱、迁徙、诉讼或官员调动的缝隙里才露出轮廓。把这种沉默看作完全顺服，正会误解隋初统一的基础。

关中资源也有地理的限制。它能供给朝廷和西北军镇，却不能同时无代价地覆盖长江、河北与塞外。北方草原的压力使长安不能把全部精力投向陈朝；南方陈朝则仍以建康为中心，掌握江面、州郡与本地官僚的联络。隋廷后来能选择南征，取决于它暂时把北门的风险控制在可承受范围内，而非因为受禅本身自动赋予了跨江能力。统一的第一层含义，是把不同方向的危险排出先后次序。

\subsection{北门没有关上：突厥、军费与南征的等待}
开皇二年，沙钵略可汗与北周旧部相连的威胁，使新朝很快看见北方边线不是一条可以写进地图就安静下来的界线。\eventanchor{ST-582-TURKIC-DEFENSE}{开皇二年的北边危机}在《隋书·突厥传》与帝纪中留下的是朝廷应对和部族关系的线索；这些汉文记录有强烈的中原朝廷视角，不能让我们直接复原草原各部的内部商议。但它们至少说明，长安必须把骑兵、粮草、使者和婚姻外交同时放在北方考量，南征陈朝因而不是一项可以孤立计算的计划。

边防的选择不是单纯的攻或守。若长期聚兵于北，关中和各道的粮运会被固定在军镇；若急于把精锐转向南方，则突厥与北周残余关系可能进入关陇腹地。隋廷使用军事防备、使节往来和对草原政治关系的分化，试图让北面不形成一个能够共同南压的联盟。史书常以某位可汗的顺逆来组织叙事，实际的压力则经过马匹、牧地、互市、部落首领的竞争和边城的补给传导。我们没有足够材料将这些因素逐村逐帐地量化，却能看出边防稳定的成本并未消失，只是暂时由不同手段分摊。

这一等待塑造了南征的节奏。长江以北的将领需要了解北线何时可以抽兵，江淮的州县必须准备船只与粮秣，陈朝也并非静止地等待。把589年的胜利解释成隋军数量或将领才能的单一结果，会把此前数年的战略选择删掉。真正的前提是：北方边防在特定时段没有把关中完全拖住，隋廷才得以把跨江战争变成可执行的工程。

边防也说明新朝为何如此关心制度语言。军队、地方官和使者不能只凭“新天子”的名义获得资源；他们需要可被识别的官号、赏罚和文书程序。北方的警报没有把法律问题留在后方，反而把它推到前台：怎样判处逃亡与盗窃，怎样处理军户、徭役和官吏失职，都会影响谁愿意留在可被调度的网络里。法律不是边防的对立面，而是让有限资源不至于完全由私下关系分配的一种尝试。

\subsection{开皇律令：纸上的轻刑与县廷里的摩擦}
开皇三年修定律令，\eventanchor{ST-583-KAIHUANG-CODE}{开皇律令的颁行（583）}在《隋书·刑法志》《北史·刑法志》和《通典·刑法》中留下制度史的清晰痕迹。志书告诉我们删并条目、整理刑名和继承前代法律的方向；它们不能告诉我们每一个县的狱舍当天如何办案，更不能证明一项较轻的法定刑自动减轻了所有人的恐惧。律令先是一套中央用来辨认问题、分派责任的语言，之后才可能在州县、军府和乡里的争执中获得不同程度的执行。

想象县廷时，最可靠的不是替无名当事人编出一场审讯，而是看见一份案件所必经的物件：状牒、印信、保人、簿籍、押送者和能够解释条文的吏。有人因逃亡、斗殴、田界或役使被带到官府，地方官要在新旧惯例之间选择引哪一条法；文书要沿着州郡上传，家族与邻里则可能以担保、赔偿或关系设法改变结果。法律的统一在这个层面从不是整齐的复制，而是不同人争夺何种解释能被登记为正式结论。

开皇律令之所以成为后世制度史的显著节点，也因为它为唐律及其注疏提供了可追溯的前史。然而这种后续影响不能代替隋代的实际图景。唐人保存和整理隋制时，已带着唐朝自己的官僚经验；《通典》又把不同年代的制度放在长时段叙述中。读者可以据此理解规则怎样被编排，却应把“颁行”“复申”“地方援用”和“实际纠纷的结果”分开。我们尤其不应从几条规范推断农户、妇女、流民和边地群体都以同一方式进入法律。

法律还有财政的一面。户籍、婚姻、逃亡、盗窃和劳役之所以频繁进入官府，并非它们天然属于不同的世界：一户从簿籍上消失，会改变差役与租调由谁承担；一名壮丁逃离，邻保与同里可能被追问；军镇需要粮帛时，地方官更迫切地把人固定在可计数的名册里。制度想要减少任意惩罚，地方行政却仍可能因征发而扩大控制。隋初的法制成就与这种张力同时存在，不能用“宽简”一词把后者遮去。

\begin{debate}{开皇律令能不能证明基层生活已经整齐化？}
制度志保留了条文沿革，使我们能讨论中央如何重整刑名；后来的唐律材料也显示其影响。可是规范文本不是执行台账。敦煌、吐鲁番等保存下来的文书主要来自特定地区与年代，不能直接充作583年关中或江南县廷的样本；墓志和正史又更容易保存精英的经历。因此可以说隋廷建立了一套具有持续影响的制度语言，不能说它已使全境案件、赋役与身份关系整齐划一。
\end{debate}

\subsection{越过长江以后：城门、仓门与590年的反复}
开皇八年十月，隋军以水陆并进发动灭陈战争。\eventanchor{ST-588-SOUTHERN-CAMPAIGN}{隋军南下（588）}的本纪叙事给出将领和战事大势，具体行军路线、船只数目和各地反应却在不同材料中并不等密。跨江战争的关键并不只在军阵碰撞，还在渡口能否供船，沿江州县是否能供应，江面风水如何影响会合，以及陈朝的守军、地方官和居民如何理解正在到来的北方军队。把战图画成一支自北向南的箭头，很容易抹掉这些使战争得以移动的条件。

次年正月建康易手，陈后主被俘。\eventanchor{ST-589-CHEN-CONQUEST}{建康易手（589）}是王朝更替的明确节点，却不是江南已被无缝纳入隋廷的证据。宫城中的俘获可以迅速改变谁签发诏令；城外的州县、寺院、仓储和家庭要经历的是官员替换、簿籍重编、军队驻扎与差役重新分配。旧陈官吏有的被吸纳，有的失去位置；本地精英要判断新朝的税役与礼制如何作用于既有网络。正是在这些不易被本纪完整写出的环节，军事占领才可能变为行政接收。

590年的江南反隋起事说明接收并不平滑。\eventanchor{ST-590-JIANGNAN-REBELLION}{江南起事（590）}的记载常把叛首、讨平和中央处置排成一条线，容易让人把复杂地方冲突缩成旧朝遗民的反抗。我们不能从这些记录判定每一处动乱的同一原因：赋役、官吏更替、地方武装、旧有政治联系与粮食压力在不同地区可能组合得很不一样。但起事本身足以提醒读者，589年的“统一”是军事和名义的统一；它之后还要经过对人的登记、对粮的征收、对道路和治所的控制，才能变成较稳定的统治。

江南在这一过程中也并非只承受。长江下游的水路、市场、作坊与士人网络会被新朝纳入更大的资源调度，并在后来的运河与都城建设中取得新的位置。问题是这种接入由谁决定、成本由谁先付。隋廷希望从南方得到粮帛和交通优势，地方社会则要面对更远的征发出口。统一所创造的不是一个没有差异的空间，而是一条能让差异被更快送往中枢、也让中枢压力更快回到地方的通路。

\subsection{东都的工地：水位、仓廪与被迁动的人}
大业元年营建东都洛阳，\eventanchor{ST-605-EASTERN-CAPITAL}{东都营建（605）}不是一座新宫城的孤立工程。洛阳处在关中、河北与江淮之间，选择它意味着朝廷试图把来自不同区域的粮食、官员和信息放到更易汇合的节点。建设需要城墙、宫室、道路、仓储与人口迁置；每一项都把工役、材料和组织能力从周边吸引进来。正史叙事能够确认朝廷的计划和宏大规模，考古遗址能提示城市道路、水系和分期；两类材料都不能替我们精确计算每一个工人的役期与死亡，也不能把“迁民”写成无差别的单向受害或受益。

都城工地的日常由许多短距离的安排构成。木料、砖石和粮食必须在合适季节抵达；仓门的开启需要簿记；被调来的工人要在陌生的水土中寻找宿处；负责工程的官吏要把上层期限翻成一段段可以验收的劳作。某处堤岸或道路留存到今天，并不能证明当年的劳动组织完全相同；但工程遗存使我们不必只在帝王意志里理解东都。它是一套把河流、土地、人力与文书临时结在一起的装置。

东都的意义还在于它改变了“远近”。从长安看，江淮与河北的资源需要跨越更长的陆路和山河；从洛阳看，水道和陆路可以被重新编排。缩短中央的指挥距离并不等于缩短每一个村庄的负担。相反，越能被纳入调度的地方，越可能在需要时接到更快的征夫、船只和粮帛命令。大业初年的建设把这种能力推向高处，也把它的脆弱性一并暴露：只要季节、战争或地方服从出现裂缝，网络的扩张就会使缺口传得更远。

\subsection{通济渠不是一条抽象的“南北大动脉”}
通济渠的开凿与东都经营相连。\eventanchor{ST-605-TONGJI-CANAL}{开凿通济渠（605）}把洛阳、淮河与江都方向的水系编入一条更密的调度路线。工程的价值不在于给后人一张壮阔的地图，而在于改变谷物、布帛、军人、官员和消息在何时、由谁、沿哪些岸边移动。河道必须疏浚，堤防必须修筑，船只与纤夫必须调集，沿线州县还要在汛期、枯水和征役之间不断调整。水面连通是一项物质事实，能否稳定运输却是由维护、季节和地方供给共同决定的过程。

运河考古和遗存提供了重要的空间证据：它们让我们看到河道、堤岸和节点并非只是史书里的修辞。可是遗迹的始建、修缮和使用往往属于不同阶段，不能仅凭一段河道便断言某年调集了多少劳力，或一户家庭遭受了何种损失。《隋书》的工程叙事同样带着唐人以隋亡为终点的安排。两种材料放在一起，最稳妥的读法不是把运河赞为“繁荣”或斥为“暴政”的单一象征，而是问它使哪些资源可以流动，又要求哪些人反复维持这种流动。

沿渠的社会经验不可能完全一致。仓场与渡口附近的人可能因运输、交易和驻泊而得到新的机会；被摊派船只、役夫或粮食的家庭则会感到国家以更近的方式进入日常。水路提高了军粮和行政文书的速度，也提高了在紧急时将地方资源转向远方的能力。后来的唐朝继续利用部分水运网络，说明隋的设施不是随政权覆亡而化为乌有；隋末的动荡又说明一条通道并不能自行生产信任、粮食或服从。

这正是制度得失的所在。东都和通济渠尝试解决的是统一后资源分散、都城与江淮相距过远的问题；成本以工役、地方供应和持续维护的形式落在沿线及其周边；副作用则是当中央把更大的计划压入同一网络时，断裂也会沿网络迅速扩展。读者随后在辽东远征和山东起事中看到的，不是两件与运河无关的新故事，而是调度能力被用于不同目标后出现的不同压力。

\subsection{废州置郡：名称改换以后，谁来认识一户人家}
大业三年废州改郡，\eventanchor{ST-607-COMMANDERY-REFORM}{废州改郡（607）}常被概括为行政层级的一次回摆。对中枢来说，名称、隶属和官署秩序可以写在诏令上；对地方来说，改变要经过印信的替换、文书格式的重写、官员的调任、旧有辖区的重新归类以及户籍和仓储向新机构报送。行政史容易把这些步骤压缩为一个制度动作，实际上每一步都要求有人知道村落、田地、渡口和逃亡者原先归谁管理。

地方治理从来不是中央命令的空白承载物。县令依赖的书吏熟悉地方亲属和田界，豪强与寺院掌握的土地、佃附和施舍网络也会影响谁能被登记、谁能获得担保、谁在征役时被优先找到。改称郡县并不会把这种知识收回到洛阳；它更可能迫使地方把既有关系翻译成新的上报语言。有的人借新官到任重新进入仕途，有的人借旧有乡里力量拖延或规避。正史很少为这些缓慢的翻译留下完整账目，因此不能将“废州置郡”直接等同为行政渗透突然加深。

然而名号的调整也并非无效。它给财政、司法和军事文书提供较一致的收发路径，让朝廷能把某一地区视作可比较、可征调的单位。统一国家需要这种可计算性，但可计算性总是选择性地看见人：能进入户籍与仓簿的被看见，流动者、躲避者和边缘群体则可能只能在逃亡、盗贼或叛乱的记录中突然出现。隋末民变因此不能仅被解释为秩序“突然”失去；它也暴露了此前登记和调度始终未能完全覆盖的社会。

\subsection{河西的夏天：吐谷浑战争与通向西域的走廊}
大业五年，炀帝亲至河西，隋军击吐谷浑并设置西海、河源等郡。\eventanchor{ST-609-TUYUHUN-CAMPAIGN}{河西与吐谷浑之役（609）}把隋的视线从都城和江淮转向青海、河西与西域通道。汉文正史以隋廷出兵、设置郡县的语言保存这场行动，因而特别容易显示中央如何命名胜利；它不能单独说明草原、绿洲和山地人群怎样理解新的行政安排，也不能让郡名本身变成稳定控制的证明。

河西不是一条空走廊。军队要越过的，是水源稀少、道路受季节和地形约束的地带；沿路城镇、牧地、驿传与市场既是补给可能，也是争夺对象。将边地写成中原军队地图外的白色背景，会看不见当地行动者的选择。吐谷浑各部、河西居民、商人与使者都以自己的安全和交换网络判断战争；隋廷设置行政机构，是希望把这些空间放进税役、军防和外交的可见范围，却未必拥有持续维持的人员与物资。

这次远征还联系到更广的知识与贸易世界。西域来使、译者、僧侣和商人把不同语言、宗教与货物带到河西，朝廷想要掌握道路，既为防务，也为使节往来和声望。可是朝贡词汇不能轻易读作单向统属，史书所列“来朝”也不足以描出每段道路的安全。敦煌和吐鲁番的文书、城址和多语材料能补出局部生活与交通的痕迹，但它们属于特定地点、书写者和年代；最稳妥的说法是，隋的河西经营扩大了接触和干预的尺度，未能消除边地政治的多中心性。

河西的经验与运河工程形成对照。运河试图把已在帝国腹地的资源接得更紧；河西经营则试图把距离更远、生态更脆弱的道路纳入行政与军事网络。两者都依赖转运与地方供给，也都不能由一道诏令保证。隋将国家能力投向这些工程和战事时，获得了新的通道与名义，却同时增加了需要长期维护的边缘。

\subsection{辽东的距离：城防、疾病和不能被兵额填平的缺口}
开皇十八年第一次征高句丽已显示东北远征的困难：水陆行军受风浪、疾病和补给制约而撤回。\eventanchor{ST-598-GOGURYEO-FIRST-CAMPAIGN}{第一次征高句丽（598）}不应只作为后来失败的预告。它提示隋廷已把辽东视为必须处理的战略方向，同时也暴露出北方道路与海上航行不能依照都城意图整齐配合。高句丽方面的材料保存与汉文正史不对称，故而我们不能把隋方叙述中的敌情、战果和对方意图写成双方都认可的事实。

大业八年的大规模远征把这种距离放大。\eventanchor{ST-612-GOGURYEO-INVASION}{第一次大规模征高句丽（612）}在《隋书》和《资治通鉴》中被组织为辽东城攻守、诸军推进与撤退的连续场面。辽水、城防、道路与粮秣使“出兵”不可能只是一个命令：部队何时渡河，车船是否赶得上，疾病怎样影响行列，前线又能否把后方送来的粮食保存到需要的时刻，都会改变战局。史书所报的巨大兵数和损失长期进入后世叙事，但它们具有动员宣传、失败归责和传抄夸张的风险，不应被换算成精确的人口统计。

在辽东城下，隋军面对的不是一座等待被中央地图涂色的孤城。城防将地形、储备、守军与地方社会结合起来；攻方若不能持续供应，围城的时间本身就变成对后方的压力。高句丽是有自身政治与军事组织的对手，不是隋亡故事中用来惩罚炀帝的布景。汉文材料主要记录隋军的决策，读者应同时记住这一观察位置：我们能较稳妥地说远征受制于距离、补给与防御，不能替对方社会写出无证的内部心理。

战争带回内地的后果也不必等待全军覆没才出现。征兵、转输、船只和马匹的需求会在不同地区落到不同名册上；有人被直接编入行列，有人负责运粮，有人留在家中承担失去劳力后的税役与耕作。材料很难让我们逐县统计这种负担，正因为如此，更不应把它抹成一句“民不聊生”。能够确认的是，辽东战场把中央的战略选择延伸到远离前线的道路、仓场和家庭，地方承受的差异则仍需保留在叙事中。

\subsection{长白山的歌与簿籍之外的人}
大业七年王薄在长白山起事，\eventanchor{ST-611-WANG-BO-REBELLION}{王薄起事（611）}常被后世叙事当作隋末民变的开端。史书保存了其名与活动，却不允许我们把一首流传的歌谣当作全山东农民的共同宣言。歌谣、传闻和官府报告能显示不满被怎样说出和传播，不能替代完整的地方社会调查。真正值得注意的是，起事让原先被视为逃亡、盗窃或治安案件的人群，开始以持续武装和地域联系挑战官府的分类。

长白山及山东平原的地理使流动者有藏匿、结集与转移的可能，也限制他们的供给。山林不是天然的自由空间；人要获得粮食、武器、消息和熟悉道路的向导，仍需与村落、市场、渡口和地方网络发生关系。官府若只以捕盗的方式处理，可能把逃亡者推向更稳定的结伙；若大量征调乡里自保，又会把本已紧张的劳役与粮食负担压得更重。起事在这里是国家分类和地方生计相互摩擦的结果，而不是抽象“民心”的突然翻转。

王薄之后，各地首领和队伍的来历、目标与活动范围并不相同。有些依托旧有豪强或地方官，有些与运输路线、盐场、仓储或水网相连，有些则在战乱中迅速更换依附。把这些人一律称为“农民军”或“盗贼”，都会抹掉他们与州县、军府和家庭之间复杂而短暂的关系。正史的秩序语言会把他们纳入平叛叙事，读者应在这种语言之外追问：哪些人不再能靠簿籍、差役和市场维持生活，哪些通道让他们能聚集，又有哪些地方选择与他们交易、躲避或对抗。

\subsection{杨玄感在黎阳：仓粮、漕路与一场没有等到辽东结局的叛乱}
613年，炀帝再次出征高句丽，杨玄感趁机在黎阳起兵。\eventanchor{ST-613-YANG-XUANGAN}{杨玄感之乱（613）}之所以震动朝廷，不只因为一位高门将领反叛，还因为黎阳靠近仓储、河道与北方交通的关键节点。把粮食、船运和军队接在一起的网络本来服务远征，反而让反叛者有机会以同一网络威胁中枢。叛乱很快被平定，却显示大规模动员并不会只把力量集中在皇帝手中；掌握节点的人也能把资源转向另一种政治选择。

杨玄感的动机在史书中被不同叙事塑形：家族记忆、个人野心、对炀帝政策的判断和临战机会都被后人放入解释。我们没有理由把其中某一项当作无可争辩的内心真相。更可观察的是他选择的时点和地点。主力在东北，内地仍须通过漕路与仓廪维持转输；朝廷在接到消息后要决定哪些军队回援、哪些水路必须守住、哪些地方官可以信赖。政治危机于是显出物流的形状。

这场叛乱也改变地方人的风险计算。沿线官吏要选边或等待，运输者和仓吏的物资可能被征用，村落在官军与叛军之间不得不提供信息或粮食。史书很少完整记录这些无名选择，但不能因此把黎阳仅看成贵族政治的插曲。它把远征、漕运和地方控制放在同一时刻，说明隋廷所谓的统一调度，是靠许多并不稳固的中间人维系的。

\begin{sourcenote}{如何同时使用《隋书》、编年史与河道遗存}
《隋书》帝纪和列传、\textit{资治通鉴}的613年条，能够确定杨玄感起兵与辽东战争相叠的基本顺序；它们尤其擅长记录朝廷如何把事件看作叛乱。仓城、河道和交通遗存能帮助定位资源为何具有战略意义，却不能证明某座遗址在每一日为哪一军提供了多少粮食。因而本节把黎阳理解为一处资源节点，而不把叙事中出现的兵粮数字写成精确统计。
\end{sourcenote}

\subsection{第三次远征的撤回：胜利语言如何掩盖一条疲惫的路}
大业十年第三次征高句丽，\eventanchor{ST-614-GOGURYEO-THIRD}{第三次征高句丽（614）}以对方遣使送还斛斯政、隋军撤回收束。朝廷可以把使者到来写成外交成果，远征者则要沿着已经多次使用的道路、渡口和补给线返回。两种事实并不互相取消：外交接触确实改变了停止作战的条件，但它不能抹掉此前征发留下的疲惫，也不能证明辽东问题已经解决。编年史在此提供年月次序，不能替我们推断将士、役夫和边地居民对撤军各自有何感受。

连续远征的负担特别难以用一个数字表现。一次征役可能意味着壮丁离开田地，另一次则可能意味着同一家庭交出车辆、布帛、粮食或接待经过的队伍；村落的损失未必在同一年显现，也未必由同一户承担。正史中的户口和军额往往属于行政统计或政治叙述，口径未必稳定。与其把不可靠的巨大数目当作感情的替身，不如看见反复动员怎样把国家与家庭之间本就脆弱的协商推得更紧：名册必须更新，逃亡更受追查，留下的人要补足空出的劳力。

撤军还改变了官员与将领的处境。有人要为战果解释，有人要看守回归的物资，有人要在沿途处理疾病、逃散与纠纷。中枢可以发布新的赏罚，不能立刻恢复每一处市场、农时和地方信任。隋末的危机不是一条从辽东直接流向长安的直线；它在不同地区经由运输、征役、地方竞争和传闻取得不同形状。第三次远征结束后仍能看见这种未被消化的压力，正说明“撤回”不是回到原状。

\subsection{雁门的围困：边境关系进入皇帝所在的城门}
大业十一年八月，始毕可汗围炀帝于雁门。\eventanchor{ST-615-YANMEN-SIEGE}{雁门被围（615）}把边防风险带到皇帝亲身所在的城池。此前朝廷总能把突厥、河西或辽东安排为远方事务；围城迫使它在很短时间内依赖守城者、勤王军、使者和外交安排。城门、粮食和消息的速度在这里比宏大的帝国名义更直接地决定局势。史书所述的解围过程主要从隋廷立场出发，不能据此简单描画草原政治的统一意图，但围困本身足以表明北方关系远非已经被稳定处置。

雁门是道路交会处，也是边地社会。城内有官员、军士、商旅与居民，城外的牧地、关口和行军路径连着不同群体的生计。围城时，皇帝的安全与普通人的粮食、守备和逃离机会被放入同一空间；不过不同人的资源并不相同。能进入朝廷通信网络的人有机会影响求援，依赖市场和日常耕牧的人则更容易被战事切断。我们不能用一条围城记录为所有人写出同一种恐慌，却可看出边疆不是帝国边缘的静态线条，而是国家能力随时可能被检验的生活地带。

雁门之后，勤王和解围并未自动修复中央权威。需要被动员的将领会记住朝廷在危急时给出的承诺，地方势力会评估谁真正拥有兵粮，突厥方面也会根据联盟与利益重新选择。危机的制度得失很清楚：集中皇权能在短时间发出广泛求援，却必须依靠地方军政网络真的肯来；外交可以争取喘息，却不能将边境关系永久固定。隋末各地的分裂，正从这些被迫外露的依赖中继续生长。

\subsection{寺院、施舍与国家看不见的日常账本}
隋初至大业间的宗教生活不能只由皇帝对佛、道的态度来概括。寺院是诵经、造像、安葬、施舍和旅行者停泊的场所，也可能拥有田地、劳力、施主与地方声望。造像题记、石窟、墓志和后出的僧传能让我们瞥见某些捐献者、僧人和家族；它们保存得更多的是有能力留下文字和物件的人，不能据一方题记便说整个地区信仰同样热烈，亦不能把神异叙事当作逐日记录。

对地方家庭而言，寺院网络有时提供的不是抽象教义，而是可被使用的关系：为亡者作功德、借助识字者书写愿文、在迁徙或饥荒中寻求救济、通过共同供养建立彼此可见的名声。国家的户籍和赋役簿很难完整捕捉这些往来；但当朝廷需要清查人口、控制资源或重新安排寺院时，宗教空间又会进入行政视野。由此可见，社会并不只由官府与反叛者两端构成，还由许多能够暂时承担互助、记忆和物资周转的团体连接。

隋代的材料尤其要求克制。后来唐人整理前朝兴亡时，容易把宗教政策放入贤君、昏君的道德框架；僧传又以宗教共同体的典范为重点。我们可以说佛道机构参与了隋代社会的物质与象征生活，不能从这些文本推定它们在所有县份拥有同样的财富或政治影响。把这一限度保留在叙事中，反而能看清为什么国家越希望把资源集中到可见账本，越会遇到不完全属于官署的地方关系。

\subsection{一份户籍没有留下来的地方：妇女、迁徙者与被统计的家庭}
隋代正史的中心是皇帝、将相和战场，地方社会更多以赋役、灾荒、盗贼或“户口”的形式出现。一个“户”看似稳定，却可能包含年龄不同的男女、亲属、雇工、寄居者与暂时离散者；户籍把他们编成行政单位，并不能解释家中谁实际承担耕作、照料、织造、运输或应对征役。制度史可以帮助我们理解国家希望如何登记家庭，却不能取代这些家庭的自述。

妇女在本纪中常只在皇后、宫廷婚姻或烈女故事里显影，墓志、造像题记和契约类材料则偶尔留下更具体的关系：婚姻、丧葬、施舍、财产和迁居。不过这些材料有精英和保存偏差，不能把某位贵族妇人的资源推广为乡村妇女的处境。更应避免在材料沉默处虚构“普通人”的感受。可以确定的是，大规模征役与战争不会只影响被编入军籍的男子；留守家庭对田地、债务、儿童和老人的安排，构成国家动员能够持续或断裂的隐秘条件。

迁徙者同样难以被稳定看见。有人因官职、军役和工程离乡，有人因饥荒、冲突或逃避差役而移动；官府往往在他们离开原籍、未入新籍时才以逃户或流民的名称记录。这样的标签是治理的视角，不是对个人身份的完整描述。隋末各地队伍之所以能吸收不同人群，正因为许多人已经处在原有名册和保护关系松动的状态。地方社会不是等到618年才突然破碎，它在此前的迁动、补役和重新结保中已不断调整。

\begin{textwindow}{“户”不是一个天然完整的社会单位}
《隋书》食货、刑法等志书保存国家处理户籍、赋役与刑名的术语；它们让读者看到治理者如何把家庭变成可征、可罚、可统计的单位。墓志、题记和少量出土文书能补充某些个体关系，但地区和身份极不均衡。因而本卷使用“户”时指行政上被登记的单位，不把它自动等同为共享利益、共同劳动或具有同样声音的家庭。
\end{textwindow}

\subsection{江都的断裂：军营、船队与最后一座仍在运转的城市}
大业十四年三月，宇文化及等在江都发动兵变，杀炀帝。\eventanchor{ST-618-JIANGDU-COUP}{江都兵变（618）}是隋朝皇帝生命的终点，却不是所有隋制在同一天停止。江都是运河与江淮网络的重要节点，城内有宫廷、军队、官吏、船队、市场和被征调而来的各类人。兵变发生时，这些人面对的不是一个抽象王朝的崩塌，而是粮食由谁发放、武器由谁掌握、道路能否离开、家人消息能否送达的急迫问题。

正史叙述常把兵变放在炀帝失政的结局中，容易把军士的行动简化为道德报应。军营中的不满、补给、赏罚、久役、地方隔绝与将领竞争如何交织，材料不足以让我们恢复每一个参与者的理由；更不能替他们编出密谋时的对话。但兵变能说明，长期远离关中的军政集团已形成自己的资源与恐惧，当中央不再能同时提供给养、命令和可预期的前途时，武器与船只会被用于改变权力本身。

江都之后，隋的名义还在不同人物和地区手中被短暂保存，关中、河北、江南与北方各有不同的军事政治组合。唐的兴起并非从一片真空中创造秩序，而是接收隋留下的道路、都城、法制语言、官僚经验和未解决的征役压力。隋的统一遗产因此是矛盾的：它使更大范围的资源能够被连接，也使远距离动员的代价更加清楚地压在地方网络上。对后来读者最重要的不是为隋贴上“短命”或“暴政”的标签，而是看见一个国家怎样因试图解决真实的统一、边防和运输问题，反而把自己的承重处推向极限。

\subsection{粮仓并不只在史书的数字里：丰年、储备与征调的时序}
隋初的统一常与富庶、仓储和轻徭的形象相连。这样的形象并非全无依据：中央能够积聚粮帛、修整制度和组织南征，必须以相当的财政与地方生产为条件。但“有仓”不是对所有地区、所有年份和所有家庭都同样有利的事实。仓储首先是一种将粮食从收获季节、地方市场和私人手中转入官府保管的能力；它要依赖运输、保管、簿记、官吏与可征的户口。粮一旦入仓，何时开仓、向谁发放、为军费还是赈济使用，又会决定国家储备在地方眼中意味着保障还是新的抽取。

《隋书》食货志与帝纪保存朝廷对于仓廪、赋役和户口的叙述，制度史可帮助辨认租调、力役与储粮的术语。可是这些材料的数字来自不同记录环节，常服务于评价治绩或说明危机；它们不能无缝拼成全国生活水准的曲线。一个报称充盈的中央仓也不能证明沿江、河北、关陇和边郡都在同一时刻免于歉收。反过来，某次灾害与赈济诏令也不能直接推出全境饥荒。读者要把粮仓看作连接中枢与地方的制度装置，而不是一枚可以证明“盛世”或“暴政”的总括标签。

粮食到达军队前经历的环节尤其容易被省略。县里征得的粟米要集中，仓吏要验收，车辆或船只要把它运到较大的节点，再由沿途军政机构接手；损耗、延误、盗窃、雨水和路况都会让纸面数额变形。负责运输的人可能来自军户、差役或市场雇佣，材料很少完整区分。远征辽东时，后方的粮道比战场上的命令更能检验这种体系；江都兵变时，粮食与赏赐又成为军营权力的直接对象。所谓财政能力，因而不是一个留在国库里的抽象量，而是不断被搬动、记录、争夺和消耗的物。

隋的制度选择试图以较统一的名册和征收方式让粮食可预测地进入国家。它解决了分裂时期区域资源难以共同调度的问题，也把地方对丰歉和劳力的缓冲压缩到更大的计划中。丰收时，集中可能显得有效；战争、工程和灾害相叠时，原本为安全而设的簿籍会使逃亡者和欠役者更清晰地暴露。制度得失并不取决于某一条诏令的善恶，而在于国家能否让征收、储备和救济的节奏不至于同时失去弹性。

\subsection{道路上的国家：驿传、渡口和一封没有保存下来的文书}
隋的统一依赖道路，然而道路不是把地名连成线那么简单。使者从长安到州郡，要有可换的马匹、食宿、文书和能够辨认印信的人；军队经过渡口，要有船、绳索、熟悉水势者和预先筹集的粮草；商旅带着布帛、盐、器物和消息，也会借同一套路径移动。驿站、关隘和桥梁使中央命令看起来能够迅速抵达，但它们的运转又取决于地方是否有余力提供人马和供给。

一封没有保存下来的文书可以帮助我们理解这种限制。它或许从东都发出，写明某郡须在期限内报送粮船；到达郡府后，要有人验看封缄、抄录命令，计算本地仓储、河水和可用船只，再把负担分到下属县。县中书吏必须把抽象数量转换为具体的人、物和日期，船户、里正与役夫则要在农时、家庭与官差之间安排。我们不能把这幅过程写成史实的逐日纪录，因为具体文书不存；但正是这些一般而必需的环节，说明诏令从中枢变成行动时必然要穿过许多地方判断。

道路还生产不对称的信息。中央接到的多是能被写成报告的事件：叛乱、灾害、军情、任免和数字；沿途普通人的迟误、交易和怨言若未造成更大后果，往往不入档案。地方官知道本地河道何时浅、哪一处桥坏、哪个家族能提供船只，却未必愿意把不利消息完整上报。隋末的突然性因此有一部分来自记录的突然性：许多长期的磨损只有在运输中断、队伍聚集或官军失利后才被中央看见。

运河扩大了道路的能力，也没有取消这种不对称。水路让大批物资更易汇集，却要求更大规模的堤防、船只和沿线服务；一处节点堵塞，影响会沿更长的链条传导。唐代继承隋的道路与水运，证明这些设施具有超出一朝的价值；隋朝本身却未能让设施的维护、地方负担和战略雄心保持平衡。道路在此既是统一的骨架，也是崩解时最先传递震动的神经。

\subsection{宫廷里的继承，地方上的等待：604年以后谁在调整自己的关系}
仁寿四年七月，杨广即位，是为炀帝。\eventanchor{ST-604-YANGDI-ACCESSION}{杨广即位（604）}的年序没有太大疑问，围绕继承的评价却长期被后来的隋亡叙事覆盖。唐初史书知道大业末年的结局，容易把即位者的每一个选择安排为灾难的开端。读者应把继承先还原为当时的行政事实：新的皇帝要重新组织近臣与将领，确认既有诏令、官职和赏赐仍然有效，地方官和边将则观察哪些关系将得到优先。

继承时最重要的不是揣测一个人的性格，而是资源与信息如何重新排序。旧的东宫网络、关中军政集团、北方边将、江南官员和宗教机构都可能寻求新朝廷的接近；新皇帝也需要通过任免、巡视和工程显示自己的可调动能力。东都营建、运河开凿和边疆行动在此后被加速，不能简单归为个人奢欲或雄心：它们是试图把不同区域与权力节点重新编入中枢的政策选择。问题在于，这些选择以极高的同步性要求地方提供人力、粮食和顺从。

地方社会不会以同样速度感到继承。对远离宫城的县廷来说，最先到来的可能只是新的年号、官员更替或增加的工程命令；对江淮的船户和役夫来说，继承可能在通济渠工地才成为触手可及的事实；对北方军人来说，它可能表现为辽东计划和新的征调。把不同经验压成“炀帝即位后天下皆怨”，既不能说明为什么早期工程仍能推进，也不能说明为何后来危机会在不同地方以不同形态爆发。

继承的制度风险还在于承诺。新朝廷为得到支持会授官、赏赐、减免或发布新规，随后必须有仓储和文书将承诺落地。若承诺多于可兑现的资源，地方与军队会开始依赖自己的中介关系；若征发超过可承受的节奏，原本能帮助国家运行的中介人也可能成为离心力量。杨广即位不是隋亡的充分原因，但它是理解大业时期为何会把统一、交通和边疆计划同时推高的重要时间门槛。

\subsection{毁坏与继承：隋末不只留下废墟}
618年以后，后继政权面对的是已经被战争、兵变和地方分裂撕开的空间，但也面对一批没有消失的制度与物质遗产。长安、洛阳、运河、仓储、道路、郡县文书、法律术语和受过训练的官吏仍在；它们可以被新政权接收、改名、修补或重新分配。唐的建立因此既是军事竞争的结果，也是对隋遗产的选择性继承。只写“隋亡唐兴”会把这种接收过程抹平，仿佛新朝从空地上立刻拥有治理全国的工具。

遗产从不只是一套可用资产。隋留下的还有远征带来的债务感、工程与征役形成的怨怼、地方武装的经验、流民与逃户、边防未解的紧张，以及人们对于中央命令是否能兑现的怀疑。后继者可以借运河运粮，却必须决定怎样获得沿线支持；可以沿用法制语言，却必须让地方官在新的政治关系下执行；可以占领都城，却仍需重建仓储、市场与治安。制度的连续性与政权的断裂在同一时间存在。

这也解释了为什么隋不应被写成唐的单纯序幕。它完成了南北长期分裂后的统一，留下可供唐扩展的行政和交通条件；同时，它在短时间内将都城工程、边疆战争和资源动员叠加，使地方承重处持续受压。两种事实互为条件：没有统一和调度能力，就没有那些宏大工程；没有工程和远征的成本，也难以理解调度网络何以失去弹性。隋的十八年帝国生命并不小到只能容纳一个道德故事。

\begin{causalchain}{从再统一到瓦解：不要把相邻年份读成单一因果}
\begin{tabular}{p{.23\linewidth}p{.68\linewidth}}
结构条件 & 北周末军政资源的集中、南北长期分立后的道路与财政差异、草原及辽东边境的持续压力。\\
政策选择 & 以律令、郡县、都城和水路增加中枢调度能力，同时推进河西经营与对高句丽战争。\\
触发与放大 & 连续征役、运输负担、边防危机和地方军政网络的离心选择在不同地区相互叠加。\\
直接结果 & 山东等地起事扩展，黎阳等资源节点动摇，雁门围困与江都兵变使中央失去控制军队与道路的能力。\\
中期影响 & 唐等后继力量继承隋的城、渠、法制和官僚语言，也必须承受其留下的地方分裂与动员压力。\\
\end{tabular}
\end{causalchain}

\subsection{统一后的十年：不是空白的准备期}
589年到598年之间的隋朝，常在叙事里被压缩成南征胜利和辽东战争之间的一段平静。这种压缩会错过统一真正被做成日常的过程。建康易手后，隋廷要继续处置江南起事、安排旧陈官员、重接州县簿籍，并让北方边防、关中财政和南方水路不致互相争夺同一批资源。每一件事都不如攻城显眼，却决定一场胜利能否留在地方。行政命令要有人传递，地方报告要有人抄录，积欠的赋役和新旧官员的权限要有人辨认；这些缓慢动作构成了统一最耗时的部分。

南北长期分立留下的不只是两套宫廷，也留下不同的社会和经济节奏。江南水网密集，城镇、寺院、市场与地方家族的关系不可能照搬关中经验；北方军府与边镇的供给问题又不能用建康的财政解决。隋廷希望以统一年号、法令和官制减轻这种差异带来的协调成本，但中央格式越清楚，地方翻译的压力也越大。一个州县能否按时编造账簿，往往取决于熟悉旧例的人仍否在场，而不取决于诏书的文辞是否整齐。

这一时期的宗教和文化活动也不能只被当作政治稳定的装饰。寺院重修、造像、经卷流通和士人的仕进，连接着家庭的丧葬、施舍与名望。国家有时借宗教仪式显示统一后的秩序，有时又担心寺院土地与人丁脱离常规征管。材料让我们看见朝廷与显著寺院的关系，却较少保存乡村小庙、无名施主和流动僧人的日常。正因为可见性不均，叙事不能用都城的一次法会替代全境宗教生活。

598年的第一次对高句丽战争打破了这种准备期的静态印象。战争并非从无到有地占用资源，而是把此前积累的道路、军府、税役和边防判断置于新的压力测试。其撤回说明统一后的国家仍无法轻易跨越东北的距离；但此时的失败尚未必然导向618年。中间还有十余年不同地区、不同机构如何吸收或放大压力的过程，正需从后来的工程、改革与反叛中逐段追踪。

\subsection{市场不是官仓的影子：船户、匠人和沿线交换}
大运河与东都建设的叙事若只写征夫，便会把沿线经济写成没有主动性的受害者；若只写商贸，又会忽略征发和风险。渡口、仓场和工地附近出现的交换，既可能为船户、匠人、食物供应者和小商贩提供机会，也可能使他们更容易被征用、加派或卷入军队的需求。市场并不脱离国家，但它也不完全听命于国家：价格、季节、船只损坏、河水涨落和地方治安都会改变一批粮布能否按计划转手。

考古发现的城市道路、窑址、钱币和仓储遗存，给这种物质世界留下片段的证据。它们能告诉我们某地存在生产、储藏或消费活动，不能自动说明参与者身份、交易利润或国家控制程度。一个钱币窖藏的埋藏时间并不等同所有钱币的铸造和流通时间；一处窑址的产品也不等于朝廷已经统辖其销售。将遗物和史书并读，应当让空间变得具体，而不是让零散发现承担全国经济史的结论。

对沿线家庭来说，国家工程改变的不只是要不要服役，也包括怎样安排现有技能。熟悉水路的人可能被要求带船，善于木工的人可能被调往堤岸，能提供粮食或住宿的村落可能在交通高峰时承受更多来往。有人可能利用新的流量扩大交易，有人因同样的流量失去原有的安静生计。材料不足以替这些人分配明确的收益和损失，但足以拒绝两种过于简单的画面：一条河道既不会自动把所有人带向繁荣，也不会只留下同一种苦难。

\subsection{将领与地方：军府的命令为何总要借别人的手}
隋的军事行动依赖将领，然而将领从不单独拥有一支脱离社会的军队。南征需要地方提供舟楫和引导，河西行动需要熟悉水草与道路的中介，辽东远征需要仓储、车辆、马匹和沿途治安。军府可以下令，真正把命令化为行动的是一串掌握不同资源的人：州县官、乡里头目、船户、仓吏、工匠、翻译者与临时被征发者。军事史若只写名将和战役，会把这条链条切断。

这种依赖给地方带来谈判空间，也带来危险。一个地方官若能及时供给，就可能得到奖赏或更大权力；若无法完成定额，则可能被追责，转而加压于下属。军队驻扎处，市场与农田会被消耗，居民也可能借军营寻求保护、交易或逃避原有关系。没有材料可以让我们一概断言军队必然掠夺或地方必然支持；但能够确定的是，战争让日常资源的使用权反复变动，地方社会不得不在不同武装和官署之间调整。

隋末地方军事化尤其说明这种关系的后果。中央越依赖临时调动和地方协助，越难在危机时辨认谁只是奉命、谁已形成独立网络。李渊在太原起兵、各地群雄控制运输节点，都不是凭空出现；他们利用的是此前国家为了边防和远征而经营的军政资源。这里不能把所有将领预先写成野心家，而应看到制度如何在压力下把更多武器、粮食和人际忠诚交给中间层。

\subsection{灾异、传闻与政治判断：人们怎样知道局势正在变化}
隋代本纪记录日食、洪涝、旱灾和异象，传统史学常把它们放在治乱的道德框架内。对现代读者而言，天象记录可以帮助核对某些日期，却不能直接证明天意或成为政治崩溃的自然原因。气候和灾害会影响收成、道路和疾病，具体后果仍取决于仓储、救济、交通和地方组织；同一场雨水在不同地区可能是灾难、延误或普通季节变化。

更值得注意的是灾异怎样进入信息网络。朝廷收到的奏报会选择哪些损失可被量化，地方官可能借灾情请求减免，民间传闻则会把工程、战争和不安同异常天象连在一起。我们很难恢复传闻的原话与传播范围，不应凭后出的记载虚构“天下皆言”。但是传闻作为一种政治条件确实存在：当人们不再相信官府能解释或缓解危机时，同一则消息可以加快逃亡、囤积、结社或投靠武装。

史书以灾异为帝王失德的征兆，是一种叙事选择。它让后来读者容易把复杂的财政、边防和地方冲突归为一条道德曲线。保留材料的修辞性质并不会否认灾害的物质影响，反而使我们追问更实际的问题：哪条道路被冲断，哪座仓是否开门，谁能够迁移，谁只能留在原地等待。隋末的崩解在这些不均等的应对能力中发生，而不是由天空替人作出决定。

\subsection{太原之前的北方：关中、河北和边镇并非同一片腹地}
617年李渊自太原起兵，\eventanchor{ST-617-LI-YUAN-TAIYUAN}{太原起兵（617）}通常被看作唐朝的开端。但对隋史而言，太原首先提醒我们北方内部并不一致。关中靠近旧周隋的政治核心，河北连着漕运、仓储与山东动乱，太原则处在山地、北方边防和通往关中的路线交汇处。它们接受同一朝廷的年号，并不意味着面对相同的资源、风险和消息。

李渊能够调动的军政关系来自隋制，也来自地方和边镇长期形成的信任。起兵的具体谋划、参与者动机和兵力数字在后来的唐朝叙事中带有强烈的合法化色彩，不能逐项当作中性事实。但太原在此时成为政治节点并不神秘：它有军队、道路和与北方各方联系的条件，中央又因辽东战争、山东民变和江都方向的困境而难以同时牢牢掌握这些条件。起兵是多重压力交汇后的选择，不是一个孤立英雄故事。

太原的局势也迫使当地居民作出比“支持或反对”复杂得多的判断。有人要保住家产与市场，有人担心军队过境，有人根据亲属与官职寻求保护，有人则试图在消息未明时保持沉默。正史较少保存这些人，因而不能用想象填满；但任何军队的前进都要经过城门、粮食和道路，地方选择便会以供应、引导、躲避或抵抗的形式影响结果。唐的起点由此仍属于隋所创造又失去控制的政治地理。

\subsection{渠堤上的人：一项水利工程怎样穿过村庄的日历}
通济渠的工地不能只从洛阳的规划图看。河道要占用的，是已有田畴、旧水道、渡口和堤岸；工程队到来时，地方首先要面对的不是“大运河”这个后来的名称，而是地界如何勘定、土从何处取、谁负责看守工具、何时可以离开春耕或秋收。中央诏令所规定的期限把这些不同的日历压到同一时段，县吏必须把期限翻成夫役、粮食和车辆。没有一份保存完整的沿渠劳役名册让我们复原每一户的命运，正因为如此，才不能把“百万丁夫”之类的传世数字当作精确的社会统计。

土方工程也不是单纯的肉体消耗。挖出的土要运到堤上，堤要按水势夯实，河床与旧河相接处要防止漫溢；工人需要铁器、木器、绳索、食物和临时居处。熟悉本地水流的人可能被征作向导，能驾船的人可能被要求试航，邻近的市场则被卷入粮食、草料和器具的供给。朝廷以工程获得更长的运输线，地方却以许多短距离的劳作把它接起来。运河今天可见的遗存证明这一劳动改变了地景，却无法告诉我们某个工段上的所有人是谁，也不能让后人凭一段堤岸判断每一种劳役都同样严酷。

渠堤的维护比开凿更能说明国家与地方的关系。洪水、淤塞和船只往来都会让河道不断需要修补；一项在开工时被称为国家工程的事，随后要靠沿线官府、船户和居民反复照料。若维护的责任没有稳定的资源，原先方便运输的水道也会变成新的负担。隋亡后唐仍利用水运网络，并非证明隋代征发无关紧要，而是说明基础设施可以在不同政权下被重新分配、重新解释。后继者继承的是渠道，也继承了谁应为渠道付出的难题。

工程现场还有一个常被忽略的时间问题。上层将成渠看作一个完成时刻，参与者却在不同时间到达、离开或被迫留下。有的家庭只在一季承担运输，有的人因军役或工役长期离家，有的村落因为渡口和仓场被固定在更频繁的往来中。我们无法为无名者虚构感受，却可以避免把他们并成同一张“民夫”面孔。对一个被征去挖土的人，工程是离开家庭的时日；对沿渠卖粮的人，工程可能是价格与机会改变；对县令，它是必须在限期内完成、又怕地方逃散的考核。相同河道由不同位置的人经历。

\subsection{粮船、浅滩与账簿：运输并不在河面上自动发生}
水路运输的难处常在船离开仓门之前。粮食需要脱粒、收储、核验，布帛需要按规定折算，官吏要决定哪些物资留作本地开支、哪些可以发往远方。仓簿把谷物写成数目，船上装载的却是受天气、容器、损耗和人手约束的实物。若一处仓场迟延，后面等待的船夫、纤夫和军队都会被迫改变安排；若水位过低，船必须减载或转运，纸面上的一条水路就分成许多耗时的陆水接力。

这类摩擦不应被误写为隋朝独有的无能。任何前近代大规模运输都面对季节与地形，隋的特别之处在于它把运输规模、东都工程和边疆战争更紧地叠在一起。洛阳、江都、河北仓储和北方战场之间的联系越密，中央越能在顺利时集中资源；但一个节点的缺粮、疫病或叛乱也越容易沿同一网络放大。运输能力的提高不是一条直线，它让可调度的资源增加，也让对调度失败的依赖增加。

沿河的船户并非只是被动的工具。他们掌握水势、泊处和船只修理的经验，可能受雇、被征、与官府协商，或在局势不稳时选择避开某段航道。现存材料不能让我们列出每一位船户的身份和报酬；把他们想成完全自由的商人或完全无声的役夫都不准确。更可靠的说法是：国家的漕运计划必须借助这些局部知识，而这种借助给地方留下了拖延、交易和重新选择的空间。

粮船也带着消息。船员和过客知道哪一处征役加重、哪一座城门关闭、哪支军队正向何处移动；这些消息经过市场、驿站和家族网络时会被夸大、删减或重组。中央看到的是经过奏报格式整理的危机，地方则在流言和实际物价中感受局势。隋末许多地区的选择并不等到一道明确的“天下大乱”公告才开始：当船不再按时到、军队开始索粮、道路出现陌生武装时，家庭已经要决定藏粮、迁避还是投靠。

\subsection{辽水以东的对手：高句丽不是隋朝失败的道具}
高句丽战争的叙事必须从对手的城防和政治组织开始，而不是从炀帝的结局开始。辽东诸城处在河流、山地和道路的组合中，守城者能够利用熟悉地形、储备与城墙延长攻方的补给线。隋军面对的是有自身朝廷、军队与地区网络的政权，其决策不能只作为中原皇帝判断错误的注脚。汉文正史主要记录隋廷的动员和挫折，因而对高句丽内部讨论、地方社会和伤亡尤其沉默；读者必须把这种材料不对称留在视野中。

612年诸军渡辽水、围攻城池与后续撤退的经过，在《隋书》和《资治通鉴》中被编为清晰的军事连续。清晰的叙事不等于清晰的战场。渡河要看水势和舟桥，围城要看粮食、箭矢、疾病与攻具，军队分路则使彼此能否通信成为决定性问题。史书对兵员、斩获与死亡的巨大数字带有报功、归责和后出的叙事压力；即使数字不可直接采用，距离和补给的约束仍是真实的。辽东不是一个能由中央命令跨过的空白，而是迫使隋军把后方资源持续送到边缘的地理。

海上行动同样提醒我们不要把战争理解为陆路大军的附属。风向、港口、船只质量和海岸补给决定水军能够在何时会合；一次风浪造成的失散，可能使原本在纸面上配合的部队无法形成合力。汉文材料对航海细节的保存有限，不能替我们画出每艘船的航线；但它足以说明自然条件不是故事背景，而是行动者必须计算的成本。隋朝将南方水路和北方军事结合起来的雄心，正是在海河交接处遇到难以由官僚命令消除的限制。

高句丽方面的抵抗也不应被浪漫化成自然胜利。城防、地方供给和政治组织同样需要付出代价，战争会压迫辽东及其周边居民。材料不足以让我们替他们叙述统一的体验，更不应把对方社会写成铁板一块。能做的是将战争留在双边而非单边的空间中：隋试图以一次次远征改变边界关系，高句丽以自身防御和外交选择回应，双方的普通人都在文献较少记录的运输、饥病与迁避中承受后果。

\subsection{从辽东退下来的队伍：伤病、逃亡与村里的空缺}
一支军队撤回时，并不会把战争留在边境。伤病者、逃散者、驮畜、损坏的器械和未能按期归来的劳力，都会沿道路回到不同的州县。正史以战事胜败、将领赏罚和皇帝诏令组织叙述，较少记录一座村庄如何接待归者、怎样补回荒废田地或如何处理债务。材料沉默不允许我们制造普遍悲剧的细节，却要求我们承认军事动员的社会后果不是只发生在战场。

逃亡尤其揭示行政名册和生计之间的裂缝。对官府而言，逃者可能是须追回的丁壮、欠役的户或潜在的盗贼；对家庭而言，离开可能是逃避饥饿、军役或惩罚的选择，也可能是被战争打散后的无奈。逃亡者到达新地后未必立即加入叛军，他们可能投靠亲属、进入市场、依附豪强或继续流动。直到官府把他们登记为问题，他们的移动才进入史书。隋末武装能够吸收不同来历的人，便与这种长期的名册失效有关。

地方官在这种情形中面对两难。严格追查可以补足劳役和税额，却可能迫使更多人离开；放松追查则会使上级认为其失职。县吏、里正和邻保被夹在中间，要辨认谁仍在本地、谁可担保、谁已无法找到。法律和户籍在此不再是抽象制度，而是不同人争夺责任归属的语言。我们不必把每一位地方官描成压迫者，也不能把制度写成自动仁政；关键在于危机让原本可协商的关系转化为紧迫的清查和补役。

\subsection{江淮与山东：两种水网，两种危机节奏}
隋末的叛乱常被一张全国地图压平，仿佛同一波浪同时淹没所有地区。江淮和山东都与水路、运输和征发密切相关，危机的节奏却并不相同。江淮靠近江都和运河工程，长期承受都城、船运和驻军的需求；山东连接河北仓储、东行道路和地方武装的活动空间。水网提供移动与交换的机会，也提供逃避、聚集和切断供应的路径。国家想利用水道把资源导向中枢，反叛者与地方社会也会利用相同水道改变信息和物资的流向。

王薄活动于山东，并不意味着所有山东人有同一目标。长白山一带的地理、村落联系和官府追捕构成其活动条件；盐业、市场、渡口与其他队伍的关系又不断改变。将起事者统一称作盗贼，是官府维护秩序的语言；把他们一律写作饥民，也会忽视首领、地方豪强、逃亡者和商旅可能有不同利益。读者需要看到的是，州县不能再像过去那样确定地把人、粮和消息纳入自己的记录，政治才开始取得新的可能。

江淮的局势则与江都驻军和皇帝行幸紧密相连。运河把这里接到洛阳和北方，也使大量资源和人员长期滞留。当地居民未必直接参与兵变或反叛，却要在军队采购、官差、市场波动和道路危险中调整生活。隋炀帝在江都的存在一度让这里成为中央的替代核心；当军营失控时，同一地理位置又使消息和恐惧沿水道扩散。相同的交通便利在不同政治条件下呈现为机会、征调或逃路。

\subsection{群雄并起时的接收：城门并不只为胜者开启}
隋末各地的竞争者接管城池时，最先需要处理的往往不是称帝，而是守住粮仓、安排守门者、安抚或约束旧吏、让市场继续开门。城门一开，新的旗号可以进入；簿籍、债务、囚犯和军粮却不会自动改换主人。不同群雄之所以能短暂扩张，部分因为他们能够利用隋所留下的道路、仓储和官僚知识；他们又必须面对同样的问题：如何让被征的人相信继续交粮比逃走或投向别处更可承受。

这类接收的材料多由后来胜出的政权整理。唐初叙事会突出李渊及其盟友的合法性，把其他势力写成通往统一的障碍；败者留下的文书较少，地方居民的判断更难保存。不能因此把所有非唐力量化为混乱的代名词。每一支队伍都嵌入地方资源、婚姻关系、市场和军事路线；有些地区选择短暂合作，有些被反复掠夺，有些保持观望。隋的崩解不是一个中心被抽走后自然形成的真空，而是许多地方在不确定中重组权威的过程。

唐对隋的继承由此有双重含义。它接收了能够统一资源的工具，也必须从隋末经验中学习不要把所有压力同时压在同一批道路、家庭和军政中介上。后来的唐初改革并不等于马上解决这些难题，然而隋的失败让后继者清楚看见：城门、河道和法令只有在地方关系仍可承重时才有力量。读隋末群雄，不是离开隋史，而是看隋建立的统一如何在不同地点被接收、改用和反抗。

\subsection{渠工到县门：征发名单如何改变一条乡里的关系}
运河劳役落到地方时，朝廷看见的是州县应交的夫数，乡里看见的却是一张必须有人具名的名单。县令不能亲自认识所有壮丁，只能依靠胥吏、里正和保人，把年龄、籍贯、田地与既往差役翻成可报送的数目。名单由此既是国家看见人口的工具，也成了乡里重新分配风险的场所：谁家有老病，谁家刚失去劳力，谁能以财物代役，谁与地方头目有关系，都会影响名字是否被写入。材料没有保存某个村庄的完整册籍，不能替无名者断言具体遭遇；但正因名单必须由本地关系完成，工程负担不可能均匀落到每一户。

役夫离开后，田地和家庭并不会暂停。播种、收获、修屋、照料老人和儿童以及缴纳原有税役，仍需由留下的人安排。妇女、老人、少年与未被征调的亲属可能承担更多劳动，然而官修史书极少从他们的角度记录。后出的墓志偶尔保存家族迁徙和丧葬，造像题记偶尔留下施主姓名，却都过度代表有文字与财力的人。将沉默直接改写成一幅统一的“民不聊生”画面并不诚实；更可取的是承认，国家把劳力抽离家庭的选择，必然改变了家庭内部无法完全被史书看见的分工。

地方官也未必只是命令的传声筒。期限、河段和上级考核都使他们承压；若不能按期交人，可能面临追责，若过度征发又可能造成逃亡和治安恶化。有的官员会设法在不同里之间调剂，有的会依赖强宗和富户先完成定额，有的则把压力继续向下推。没有材料可将这些策略逐县复原，但“国家—百姓”的二分法会抹去实际实施者。隋的工程能力正来自这些中间层的协作，而危机发生时，最先松动的也常是他们维持的信任。

\subsection{堤岸之外的地方回应：躲役、代役与不合作的边界}
地方社会对征发并非只有公开叛乱和完全服从两种选择。有人可能逃往亲属所在的县份，有人以钱物或人情寻找代役者，有人拖延报到，有人把船只藏到支流，有人以疾病、丧事或道路不通请求宽限。这些行为很少完整进入正史，通常只有在积欠、追捕或起事时才被官府记录为问题。它们却构成国家动员的日常边界：每一次小小的延误都要求县吏重新找人、调粮或改期，积累起来便会使纸面工程与实际进度分离。

不合作也不等于所有人都反对运河。沿线的一些人依靠渡口、市场和运输获得新机会，甚至可能愿意维护水道；同一家庭在不同季节也会作出不同选择，收成好时能够应役，歉收或有人病亡时则更可能逃避。社会回应因此不能被压成一条“民众反抗”的曲线。隋廷面对的是会随粮价、水位、军情和家庭状况变化的地方计算，而不是一群永远同声的臣民。

官府的应对同样有成本。加紧追捕会使更多人离开原籍，减免或宽限又可能使其他征发无法完成；倚重地方豪强能迅速补足缺额，却让后者在乡里获得更大的支配资源。隋末许多队伍能在地方站稳，正说明此前的征发、保甲和供给关系已让一些人掌握了超出常规的武装与人际联系。工程本身不是叛乱的唯一原因，但它使国家更频繁地触及地方社会最脆弱的劳力安排。

\subsection{守仓的人：粮食在战场以前已经是一场政治}
辽东战争的补给困难，不应从大军抵达前线才开始叙述。粮食首先由各地仓储和征收体系集中，守仓者要核验数量、安排出纳、应对上级催促与地方求赈之间的冲突。仓门一旦为军粮打开，留给地方市场和救济的余量便可能减少；若仓门不开，前线运输又会受阻。这样的选择未必总以明确的诏令出现，常常由官吏在日期、账目和实际库存之间作出。

史书爱记录“积粟如山”或军粮不足的结论，较少记录粮从谁家手中进入仓、在何处损耗、谁因催粮而负债。数字的精确外观容易制造一种错觉，仿佛只要加总全国粮额就能理解战争能力。实际上，谷物怕潮，布帛怕损，车辆与船只要修理，人和牲畜要吃食；从一处仓运出并不等于在另一处被军队收到。每一段转运都可能有损耗、截留或延误，中央必须通过层层文书去追踪，而文书又可能把责任推向最弱的一环。

守仓者的处境也说明地方政治并非只由大人物决定。仓吏、运夫、保人和县令掌握不同的钥匙、印信和消息；他们未必能改变战争目标，却能影响物资何时离开、谁得到口粮、哪一项亏空被写进报告。我们不能据此把小吏浪漫化为自主政治者，但也不能把他们从国家能力中删去。隋廷能发动远征，正是由于这些不显眼的岗位不断将粮食转成可移动的军需。

\subsection{边堡不是一堵墙：驻军如何依赖周围的田地、市场和道路}
辽东与北方边线上的驻军需要城堡，却不只靠城堡生存。城墙可以延缓攻击，不能自己生出粮草；戍卒需要与后方道路相连，马匹需要草料，病者需要安置，信使需要能通过的关口。边堡周围的农地、集市、河流和山路因而都是军事空间的一部分。隋朝向东北集中军力时，不仅在地图上推进边线，也把这些地方纳入更频繁的征调和监控。

边地居民的身份不应被简化为“汉人”或“外族”的固定两类。商人、农户、牧人、译者、归附者和军属可能使用不同语言、拥有跨境亲属或同时依赖多个市场。汉文史书常按朝贡、叛服和军事威胁分类他们，分类本身服务于朝廷的政治视角。我们可以利用史书定位行动和道路，不能把其中的族名和官号直接投射成稳定的现代身份，更不能据朝贡词汇断言边地已被单向统属。

驻军的家庭成本也很具体。戍卒离乡，军属可能随营或留在原籍；补给迟到时，官府与地方市场之间会发生竞争；一座边堡失守或被围，附近居民的逃避路线便被切断。可见材料不足以使我们逐堡描画生活，但足以让边防从帝纪中的一条战讯回到道路、粮草和家庭安排上。隋对高句丽的战争之所以难以维持，正因为这种依赖链越拉越长，越难靠一次命令补齐。

\subsection{远征中的家书与缺席：材料沉默怎样限制我们的叙事}
隋代没有为辽东远征留下像后世那样成体系的士兵家书。我们知道征调、战役和撤退的年月，却很少知道一名士卒离家时如何同亲属商量，或者一位留守者怎样等待消息。这个空白不应被虚构的对话填满。它本身是材料结构的一部分：官修史书关心的是皇帝、将领和战局，普通人的文字难以保存，地方文书也多零散并且地域有限。

承认空白并不等于只能沉默。军役必须从家庭中抽取劳力，运输必须占用道路和粮食，伤病与逃亡必须在原籍或新地被处理；这些制度和物质条件让我们能够谨慎地谈论“家庭成本”。这里的成本不是一个可精确计算的总数，而是劳动力缺席后由谁补耕、谁照料、谁负担税役的持续问题。不同家庭拥有的土地、亲属和资源不同，承受力也不同；正是这种不均匀，解释了为什么同一轮征发会在不同地区产生不同后果。

\subsection{江都之前的军营：赏赐、欠饷与士兵的政治位置}
长期驻扎的军队会形成自己的时间。士兵关心的不只是下一道战略命令，还有口粮是否按时、赏赐能否兑现、疾病是否有人处置、归期是否仍可信。朝廷以军功、官爵和赏赐维持忠诚，也以纪律和惩罚维持秩序；当运输受阻、战事反复、皇帝远离关中时，这些手段会同时失灵。江都兵变不能只归因于一个人的煽动，它发生在一支长期被拉出原有社会关系、又逐渐不相信中央能兑现承诺的军队之中。

《隋书》对兵变的叙述带着后继王朝的亡国判断，无法为我们提供军营中每一种不满的准确比例。不能把所有士兵都写成饥饿或所有将领都写成野心家。然而军营拥有武器、船只和集体行动的条件，这使它在中央危机时具有普通地方家庭没有的政治力量。兵变将这种力量从边防与远征的工具，转成改变最高权力的手段。隋制原本依靠军队保护统一，最后却必须面对军队自身的供给和归属问题。

\subsection{唐的接收：不是把隋的一切照单全收}
唐初接收隋的都城、运河、官僚和法律语言，并不意味着唐简单延续隋的统治方法。新政权要判断哪些仓储还能使用，哪些将领可以安置，哪些地方势力必须谈判或消灭，哪些徭役可以暂缓以恢复生产。接收首先是一连串选择，而非一份遗产清单。隋留下的设施有现实价值，但它们已经嵌在战争、逃亡和地方不信任之中；唐若要使用，便必须重新取得沿线和州县的合作。

唐初史书在叙述这一过程时有强烈的合法化目的：它倾向把唐的胜利写成对隋弊的纠正。这个框架能帮助读者看见后继者如何理解前朝，也可能遮蔽唐同样继承的征调、边防与财政难题。唐的早期措施不宜被写成一瞬间修复天下；不同地区从隋末战争中恢复的速度不同，地方社会对新官府的观察也不同。真正的接收发生在城门、仓场、县廷与道路上，发生在那些愿意继续交粮、服役或提供信息的人重新评估新政权的时刻。

因此，隋末与唐兴之间不应有一条干净的断线。隋的统一工程留下了国家可以利用的通道，也留下了人们对过度动员的记忆；唐的建立既依赖这些通道，也需要避免重演其承重处断裂的方式。将这一连续性看清，才不会把隋缩成唐的背景，或把唐写成无代价的救赎。
